Generating synthetic financial time series that faithfully preserve the statistical properties of real market data is a central challenge in quantitative finance.
High-fidelity synthetic data are essential for stress testing risk models against plausible but unobserved market scenarios, validating portfolio optimization algorithms, and augmenting limited training sets for machine learning pipelines \cite{assefa2020generating, jordon2022synthetic}.
Empirical equity excess growth rates exhibit three well-documented statistical regularities, namely heavy-tailed (leptokurtic) distributions, negligible linear autocorrelation in raw returns, and persistent volatility clustering, that any generative framework must simultaneously reproduce to be considered statistically faithful \cite{cont2001empirical, mandelbrot1963variation}.

Parametric approaches to synthetic data generation face a fundamental tension between distributional accuracy and temporal fidelity.
GARCH models \cite{engle1982autoregressive, bollerslev1986generalized} excel at reproducing the slow decay of the autocorrelation function (ACF) of absolute returns, the hallmark of volatility clustering, but their unimodal innovation distribution cannot represent the multi-regime structure of equity markets, leading to poor distributional fit.
Conversely, i.i.d.\ generators that match the marginal distribution (for example, Laplace sampling) achieve high distributional pass rates but produce flat ACF profiles by construction.
Hidden Markov models (HMMs) offer a middle path: the latent regime structure can simultaneously generate heavy-tailed mixtures and temporal persistence through the transition dynamics.

However, the application of HMMs to financial stylized facts has been constrained by an influential negative result.
\citet{ryden1998stylized} systematically evaluated Gaussian-emission HMMs with $K = 2$--$3$ states against the canonical stylized facts and concluded that while HMMs reproduce heavy tails and negligible return autocorrelation, they \emph{fail} to capture the slow decay of the ACF of absolute returns.
The mechanism is clear: with few states, the geometric sojourn-time distribution of a standard Markov chain produces exponentially fast ACF decay, insufficient to match the hyperbolic-like persistence observed empirically.
\citet{bulla2006stylized} subsequently showed that hidden semi-Markov models (HSMMs), which replace the geometric sojourn distribution with flexible alternatives, restore ACF fidelity at the cost of substantially increased model complexity.

\citet{alswaidan2026hybrid} proposed an alternative solution within the discrete HMM framework: a Poisson jump-duration mechanism that forces the model to dwell in tail states for empirically calibrated durations.
That approach achieved 91--95\% Kolmogorov--Smirnov (KS) pass rates across 424 US equities and partially reproduced volatility clustering.
However, the discrete framework introduced three sources of complexity.
First, Laplace quantile binning discretizes continuous observations and introduces quantization error.
Second, two additional hyperparameters ($\epsilon$ for jump probability and $\lambda$ for mean jump duration) require grid-search calibration.
Third, within-bin Student-$t$ emission models are estimated separately from the transition structure.

In this paper, we show that these complexities are unnecessary.
A continuous Gaussian HMM (CHMM) trained via the Baum-Welch algorithm \cite{baum1970maximization, rabiner1986introduction} with quantile-based initialization reproduces all three canonical stylized facts at moderate state resolution ($K = 6$--$13$) \emph{without} requiring jump augmentation, discretization, or separate emission fitting.
The key insight is that the negative result of \citet{ryden1998stylized} is an artifact of insufficient state resolution: at $K \geq 3$ with proper initialization, the Baum-Welch algorithm learns a transition matrix with sufficient diagonal dominance in high-volatility states to generate slow ACF decay, effectively approximating the flexible sojourn-time distributions of HSMMs through a mixture of geometric distributions.

Beyond single-asset fidelity, generating \emph{jointly plausible} multi-asset scenarios is essential for portfolio risk applications, yet it is not directly addressed by the univariate HMM literature.
A factor-model construction such as the Single Index Model (SIM) of \citet{sharpe1963simplified} propagates a market engine to a cross-section, but at the cost of imposing a rigid linear dependence structure and replacing each asset's carefully fitted marginal with the regression fit.
Copula theory \cite{sklar1959fonctions, nelsen2006introduction} provides a conceptually cleaner decomposition: a joint distribution factors uniquely into its marginals and a copula that encodes dependence, so the per-asset CHMM can be reused verbatim as the marginal while a Gaussian or Student-t copula \cite{demarta2005tcopula, mcneil2015quantitative} injects cross-sectional co-movement. We implement both constructions on a representative six-asset subset (SPY, NVDA, JNJ, JPM, AAPL, QQQ) and report a consistent ordering in which copulas outperform the SIM factor baseline both on per-asset marginal fidelity and on correlation-matrix reproduction.

The contributions of this work are:
\begin{enumerate}
    \item \textbf{Baum-Welch training without jump augmentation.}
    We train continuous Gaussian emission parameters $(\mu_k, \sigma_k)$ and the transition matrix $\mathbf{T}$ jointly from raw observations via the Baum-Welch algorithm with quantile-based initialization.
    The resulting CHMM achieves 90--94\% KS pass rates and ACF-MAE of 0.046--0.049 on absolute returns, comparable to GARCH(1,1) at ACF-MAE of 0.047, without any jump mechanism, eliminating two hyperparameters ($\epsilon$, $\lambda$) and the discretization step.

    \item \textbf{Contradiction of the Ryd\'{e}n et al.\ (1998) consensus.}
    We demonstrate that the widely cited limitation of Gaussian HMMs, the inability to reproduce volatility clustering, is a consequence of testing only $K = 2$--$3$ states.
    At moderate $K$, the learned transition matrix develops sufficient persistence in high-volatility states to achieve ACF-MAE comparable to GARCH, which we show is effectively an approximation to the HSMM solution of \citet{bulla2006stylized}.

    \item \textbf{Comprehensive four-model benchmark.}
    We evaluate the CHMM against bootstrap resampling, Gaussian and Laplace i.i.d.\ generators, and GARCH(1,1) using seven complementary metrics (KS and Anderson--Darling pass rates, excess kurtosis, ACF-MAE, Wasserstein-1 and Hellinger distances, and quantile coverage) on 1{,}000 simulated paths in both in-sample (2014--2024) and out-of-sample (2025) windows.

    \item \textbf{State resolution analysis.}
    We systematically vary $K \in \{3, 6, 9, 11, 13\}$ and show that distributional fidelity improves with $K$ while ACF-MAE remains remarkably stable, and quantile coverage is 100\% at all resolutions.

    \item \textbf{Cross-asset SIM and copula extensions.}
    We extend the univariate framework to multi-asset synthesis through two complementary cross-sectional dependence models that both preserve each asset's fitted CHMM marginal: a Single Index Model with the CHMM-driven SPY path as the market engine, and Gaussian / Student-t copulas fitted by Kendall's $\tau$ and simulated via rank reordering.
    On six representative assets, copulas substantially outperform the SIM factor baseline on per-asset in-sample KS pass rate (median 98.0\% versus 77.0\%) and on the off-diagonal correlation MAE (0.026 for Student-t, 0.030 for Gaussian, 0.079 for SIM), with the Student-t copula profile MLE selecting $\nu^* = 6$ degrees of freedom.

    \item \textbf{Volatility index generalization and regime-switching volatility for option pricing.}
    We confirm robustness across NVDA, JNJ, JPM, AAPL, and QQQ, and extend the framework to CBOE VIX and VXX log returns, demonstrating that the CHMM captures the distinct statistical properties of volatility indices and provides a natural bridge to regime-switching option pricing through the decoded per-regime volatility map.
    Benchmarking the resulting implied-volatility surface against a parametric stochastic-volatility model is the subject of a companion paper and is intentionally outside the scope of this study.
\end{enumerate}

The remainder of this paper is organized as follows.
Section~\ref{sec:related} reviews the relevant literature on stylized facts, parametric models, HMMs, cross-asset dependence, and deep generative approaches.
Section~\ref{sec:methods} describes the data, the CHMM specification, the Baum-Welch training procedure, the benchmark models, the cross-asset SIM and copula constructions, and the validation metrics.
Section~\ref{sec:results} presents the empirical results: descriptive statistics, four-model comparison, state resolution sensitivity, out-of-sample evaluation, cross-asset SIM and copula dependence, and the VIX/VXX extension.
Section~\ref{sec:discussion} interprets the findings, and Section~\ref{sec:conclusion} summarizes the contributions and outlines future directions.
